\documentclass[11pt,a4paper]{article}
\usepackage[utf8]{inputenc}
\usepackage[T1]{fontenc}
% \usepackage[document]{ragged2e}
% \usepackage[a4paper, total={5in, 8in}]{geometry}
\usepackage[english]{babel}
\usepackage{amsmath, amsthm, amssymb, amsfonts}
\usepackage{enumerate}
\usepackage{graphicx}
\usepackage{fancyhdr}
\usepackage{booktabs}
\usepackage{csvsimple}
\usepackage{longtable}
\usepackage{textcomp}
\usepackage{float}
\usepackage[colorlinks=true,breaklinks=false]{hyperref}
\usepackage{listings}
\usepackage[norelsize,ruled,vlined]{algorithm2e}
\usepackage{csquotes}
\usepackage{url}
\usepackage{ifthen}
\usepackage{keyval}
\usepackage{etoolbox}
\usepackage{etex}
\usepackage{datetime}
\usepackage{eurosym}
\usepackage{tikz}
\usepackage{wrapfig}
\usepackage{paralist}
\usepackage[backend=bibtex]{biblatex}
\usepackage{xcolor}
\usepackage{caption}
\DeclareCaptionFont{white}{\color{white}}
\DeclareCaptionFormat{listing}{\colorbox[cmyk]{0.43, 0.35, 0.35,0.01}{\parbox{\textwidth}{\hspace{15pt}#1#2#3} } }
\captionsetup[lstlisting]{format=listing, labelfont=white, textfont=white, singlelinecheck=false, margin=0pt, font={bf,footnotesize} }

\hyphenpenalty=10000
% \pretolerance=9999
% \tolerance=9999
% \emergencystretch=0pt
\addbibresource{project.bib}
% \setlength{\parindent}{2em}
\renewcommand{\baselinestretch}{1.1}
\newcommand{\scalafmt}{\texttt{scalafmt}}

% Single spacing after dots and colons
\frenchspacing
\hyphenation{scalafmt}
\hyphenation{ClangFormat}
% \hyphenation{DSLVirtualization}
% \hyphenation{ReificationTransformation}

%%%%%%%%%%%%%%%%%%%%%%%%%%%%%%%%%%%%%%%%%%%%%%%%%%%%%%%%%%%%%
%                        Setup
%%%%%%%%%%%%%%%%%%%%%%%%%%%%%%%%%%%%%%%%%%%%%%%%%%%%%%%%%%%%%

\begin{document}
\raggedright

\definecolor{dkgreen}{rgb}{0,0.6,0}
\definecolor{gray}{rgb}{0.5,0.5,0.5}
\definecolor{mauve}{rgb}{0.58,0,0.82}

\lstdefinestyle{scala}{
  language=Scala,
  aboveskip=3mm,
  basicstyle=\scriptsize\ttfamily,
  belowskip=3mm,
  breakatwhitespace=false,
  showspaces=false,
  showlines=true,
  columns=flexible,
  keepspaces=true,
  commentstyle=\color{dkgreen},
  keywordstyle=\color{blue},
  numbers=left,
  numberstyle=\tiny\color{gray},
  showstringspaces=false,
  stepnumber=1,
}


\definecolor{dkgreen}{rgb}{0,0.6,0}
\definecolor{gray}{rgb}{0.5,0.5,0.5}
\definecolor{mauve}{rgb}{0.58,0,0.82}

\lstset{language=Scala,
  aboveskip=3mm,
  basicstyle={\small\ttfamily},
  belowskip=3mm,
  breakatwhitespace=true,
  breaklines=true,
  columns=flexible,
  commentstyle=\color{dkgreen},
  keywordstyle=\color{blue},
  numbers=left,
  numberstyle=\tiny\color{gray},
  showstringspaces=false,
  stepnumber=1,
  stringstyle=\color{mauve},
  tabsize=2,
}

\usetikzlibrary{calc,trees,positioning,arrows,chains,shapes.geometric,%
    decorations.pathreplacing,decorations.pathmorphing,shapes,%
    matrix,shapes.symbols}
\tikzset{>=stealth',
  punktchain/.style={rectangle,
    rounded corners,
    % fill=black!10,
    draw=black, very thick,
    text width=12.5em,
    minimum height=2.5em,
    text centered,
    on chain},
  every join/.style={->, shorten >=2pt},
}


% Example of title page for the projects carried out within the lasec

% Simply include it in your mastex tex file:
%        % Example of title page for the projects carried out within the lasec

% Simply include it in your mastex tex file:
%        % Example of title page for the projects carried out within the lasec

% Simply include it in your mastex tex file:
%        \input{cover}


% Updated March 2006 (SP)


\newcommand{\logoepfl}[0]{
  \begin{center}
    \includegraphics[width=4cm]{logo_epfl_coul.eps}
  \end{center}
  \vspace{0.3cm}
  \hrule
}
\newcommand{\logolasec}[0]{
  \vspace{1cm}
  \hrule
  \begin{center}
    \includegraphics[width=4.5cm]{logo_lasec_coul.eps}
  \end{center}
}
\newcommand{\project}[1]{
  \begin{center}
    \large{#1}
  \end{center}
  \vspace{1cm}
}
\newcommand{\department}[1]{
  \begin{center}
    \large{#1}
  \end{center}
}
\newcommand{\supervisor}[3]{
  \begin{center}
    \begin{normalsize}{
        \bf #1}\\#2\\#3
    \end{normalsize}
  \end{center}
}
\renewcommand{\author}[1]{
  \begin{center}
    \Large{#1}
  \end{center}
  \vspace{0.5cm}
}
\renewcommand{\title}[1]{
  \vspace{3cm}
  \begin{center}
    \huge{#1}
  \end{center}
  \vspace{1.7cm}
}
\renewcommand{\date}[2]{
  \begin{center}
    \normalsize{#1 #2}
  \end{center}
  \vspace{0.5cm}
}


\thispagestyle{empty}


% begin title page
  \logoepfl

  \title{scalafmt: opinionated \\ code formatter for Scala}

  \author{Ólafur Páll Geirsson}
  \department{School of Computer and Communication Sciences}
  \project{A thesis submitted for the degree of Master of Computer Science at École polytechnique fédérale de Lausanne}

  \date{June}{2016}

  \begin{center}
    \begin{tabular}{cc}
      \begin{tabular}{p{4.0cm}}
        \supervisor{Responsible}{Prof. Martin Odersky}{EPFL / LAMP}
      \end{tabular}&
      \begin{tabular}{p{4.0cm}}
        \supervisor{Supervisor}{Eugene Burmako}{EPFL / LAMP}
      \end{tabular}
    \end{tabular}
  \end{center}

% end title page




% Updated March 2006 (SP)


\newcommand{\logoepfl}[0]{
  \begin{center}
    \includegraphics[width=4cm]{logo_epfl_coul.eps}
  \end{center}
  \vspace{0.3cm}
  \hrule
}
\newcommand{\logolasec}[0]{
  \vspace{1cm}
  \hrule
  \begin{center}
    \includegraphics[width=4.5cm]{logo_lasec_coul.eps}
  \end{center}
}
\newcommand{\project}[1]{
  \begin{center}
    \large{#1}
  \end{center}
  \vspace{1cm}
}
\newcommand{\department}[1]{
  \begin{center}
    \large{#1}
  \end{center}
}
\newcommand{\supervisor}[3]{
  \begin{center}
    \begin{normalsize}{
        \bf #1}\\#2\\#3
    \end{normalsize}
  \end{center}
}
\renewcommand{\author}[1]{
  \begin{center}
    \Large{#1}
  \end{center}
  \vspace{0.5cm}
}
\renewcommand{\title}[1]{
  \vspace{3cm}
  \begin{center}
    \huge{#1}
  \end{center}
  \vspace{1.7cm}
}
\renewcommand{\date}[2]{
  \begin{center}
    \normalsize{#1 #2}
  \end{center}
  \vspace{0.5cm}
}


\thispagestyle{empty}


% begin title page
  \logoepfl

  \title{scalafmt: opinionated \\ code formatter for Scala}

  \author{Ólafur Páll Geirsson}
  \department{School of Computer and Communication Sciences}
  \project{A thesis submitted for the degree of Master of Computer Science at École polytechnique fédérale de Lausanne}

  \date{June}{2016}

  \begin{center}
    \begin{tabular}{cc}
      \begin{tabular}{p{4.0cm}}
        \supervisor{Responsible}{Prof. Martin Odersky}{EPFL / LAMP}
      \end{tabular}&
      \begin{tabular}{p{4.0cm}}
        \supervisor{Supervisor}{Eugene Burmako}{EPFL / LAMP}
      \end{tabular}
    \end{tabular}
  \end{center}

% end title page




% Updated March 2006 (SP)


\newcommand{\logoepfl}[0]{
  \begin{center}
    \includegraphics[width=4cm]{logo_epfl_coul.eps}
  \end{center}
  \vspace{0.3cm}
  \hrule
}
\newcommand{\logolasec}[0]{
  \vspace{1cm}
  \hrule
  \begin{center}
    \includegraphics[width=4.5cm]{logo_lasec_coul.eps}
  \end{center}
}
\newcommand{\project}[1]{
  \begin{center}
    \large{#1}
  \end{center}
  \vspace{1cm}
}
\newcommand{\department}[1]{
  \begin{center}
    \large{#1}
  \end{center}
}
\newcommand{\supervisor}[3]{
  \begin{center}
    \begin{normalsize}{
        \bf #1}\\#2\\#3
    \end{normalsize}
  \end{center}
}
\renewcommand{\author}[1]{
  \begin{center}
    \Large{#1}
  \end{center}
  \vspace{0.5cm}
}
\renewcommand{\title}[1]{
  \vspace{3cm}
  \begin{center}
    \huge{#1}
  \end{center}
  \vspace{1.7cm}
}
\renewcommand{\date}[2]{
  \begin{center}
    \normalsize{#1 #2}
  \end{center}
  \vspace{0.5cm}
}


\thispagestyle{empty}


% begin title page
  \logoepfl

  \title{scalafmt: opinionated \\ code formatter for Scala}

  \author{Ólafur Páll Geirsson}
  \department{School of Computer and Communication Sciences}
  \project{A thesis submitted for the degree of Master of Computer Science at École polytechnique fédérale de Lausanne}

  \date{June}{2016}

  \begin{center}
    \begin{tabular}{cc}
      \begin{tabular}{p{4.0cm}}
        \supervisor{Responsible}{Prof. Martin Odersky}{EPFL / LAMP}
      \end{tabular}&
      \begin{tabular}{p{4.0cm}}
        \supervisor{Supervisor}{Eugene Burmako}{EPFL / LAMP}
      \end{tabular}
    \end{tabular}
  \end{center}

% end title page


\begin{abstract}
  Automatic code formatters bring a lot of benefits to software development.
  When done right, code formatters relieve the developer's attention from
  syntactic trivia and enforce a consistent coding style between developers.
  However, developing a good code formatter is still somewhat of a black art and yet,
  little research has been made towards the algorithms that power such code formatters.

  This thesis addresses the problem of developing a feature-rich code formatter for a custom programming language.
  % A critical feature of such a formatter is enforcing a column-width limit, which turns out to be a hard problem.
  Our contributions are twofold.
  First, we present a language agnostic framework that consists of core data structures, algorithms and tooling 
  that allow rapid development and testing of such a code formatter.
  Secondly, we use our model in a case study and provide a concrete implementation of a code formatter, \texttt{scalafmt}.
  We show that the framework makes


\end{abstract}

\begin{abstract}
  Code formatters bring many benefits to software development, yet they can be tricky to get right.
  This thesis addresses the problem of developing a code formatter for the Scala programming language.
  We present scalafmt, an opinionated Scala code formatter that captures many popular Scala idioms and coding styles.
  Scalafmt has been installed over 6.000 times since its initial release 3 months ago, adopted by several open-source libraries and attracted a vibrant community of external contributors.
  Our work has been limited to formatting Scala code.
  Still, this thesis presents algorithms and tooling that we believe can be of interest to the development of code formatters for other programming languages.
\end{abstract}
\tableofcontents
\lstlistoflistings
\listofalgorithms
\listoffigures
\listoftables

\setlength{\parskip}{1em}

\section{Introduction} % (fold)
\label{sec:Introduction}
The main contributions presented in this thesis are the following:
\begin{itemize}
    \item A DSL for declaring feature-rich code formatting rules.
    \item An case study where the DSL is used to format Scala programs.
\end{itemize}
Throughout the paper we assume familiarity with the basics of the Scala
Programming Language~\autocite{odersky_scala_2004}.

\newcommand{\naive}{na\"{\i}ve}
\newcommand{\dartfmt}{\texttt{dartfmt}}
\newcommand{\gofmt}{\texttt{gofmt}}
\newcommand{\rfmt}{\texttt{rfmt}}
\section{Background}
This chapter explains the necessary background to understand Scala and code formatting.
More specifically, we motivate why Scala presents an interesting challenge for code formatters.
We go into details on Scala's rich syntax and popular idioms that introduced unique challenges to the design of scalafmt.
We follow up with a history on code formatters that have been developed over the last 70 years.
We will see that although code formatters have a long history, a new tradition of optimization based formatters -- which scalafmt proudly joins -- started only recently in 2013.
% Scalafmt follows this new tradition.

\subsection{Scala the programming language}\label{sec:scala}
Scala\autocite{odersky_scala_2004} is a general purpose programming language that was first released in 2004.
Scala combines features from object-oriented and functional programming paradigms, allowing maximum code reuse and extensibility.

Scala can run on multiple platforms.
Most commonly, Scala programs compile to bytecode and run on the JVM.
With the releases of Scala.js\autocite{doeraene_scala.js_2015}, JavaScript has recently become a popular target platform for Scala developers.
Even more recently, the announcement of Scala Native\autocite{shabalin_scala_2016} shows that LLVM may become yet another viable target platform for Scala developers.

Scala is a popular programming language.
The Scala Center --- a not-for-profit organization focused on Scala open-source and education --- estimates that more than half a million developers use Scala\autocite{odersky_scala_2016}.
Large organizations such as Goldman Sachs, Twitter, IBM and Verizon rely on Scala code to run business critical applications.
The 2015 Stack Overflow Developer Survey shows that Scala is the 6th most loved technology and 4th best paying technology to work with\autocite{_stack_????}.
The popularity of Apache Spark\autocite{zaharia_spark_2014}, a cluster computing framework for large-scale data processing, has made Scala a language of choice for many developers and scientists working in big data and machine learning.

Scala is a programming language with rich syntax and many idioms.
The following chapters discuss in detail several prominent syntactic features and idioms of Scala.
Most importantly, we highlight coding patterns that encourage developers to write a single large statement over of multiple small statements.

\subsubsection{Higher order functions}
Higher order functions (HOFs) are a common concept in functional programming languages as well as mathematics.
HOFs are functions that can take other functions as arguments and return functions as values.
Languages that provide a convenient syntax to manipulate HOFs are said to make functions first-class citizens.

Functions are first-class citizens in Scala.
Consider listing~\ref{lst:hof}.
\lstinputlisting[label={lst:hof}, float, caption=Higher order functions]{code/hof.scala}
The method \texttt{twice} takes an argument \texttt{f}, which is a function from an integer to an integer.
The method returns a new function that will apply \texttt{f} twice to an integer argument.
This small example takes advantage of several syntactic conveniences provided by Scala.
For example, in line 2 the argument \texttt{\_ + 3} creates a new \texttt{Function[Int, Int]} instance.
The function call \texttt{f(x)} is in fact sugar for the method call \texttt{f.apply(x)} on the \texttt{Function[Int, Int]} instance.
Listing~\ref{lst:hof_nosugar} shows an equivalent program to listing~\ref{lst:hof}, without using syntactic conveniences.
\lstinputlisting[label={lst:hof_nosugar}, float, caption=Higher order functions expanded]{code/hof-nosugar.scala}
Observe that the body of \texttt{twice} was expressed as a single statement in line 1 of listing~\ref{lst:hof} but as two independent statements in listing~\ref{lst:hof_nosugar}.

\subsubsection{Term blocks}
Scala allows term blocks to appear anywhere in a Scala code.
A term block is a sequence of statements wrapped by curly braces \texttt{\{\}}.
Listing~\ref{lst:block} shows two examples of term blocks.
\lstinputlisting[label={lst:block}, float, caption=Term blocks]{code/block.scala}
Variables bounds inside a term block do not escape the block.
Therefore, the variable \texttt{y} can be assigned both inside the first block as well as to the return value of the \texttt{function} call.
The lightweight syntax to create term blocks in Scala make them a popular feature among Scala developers.
Observe that without term blocks, the second argument to \texttt{function} would be defined externally.
Instead, the second argument is defined inline making the entire \texttt{function} call significantly bigger.

% Blocks are .
% \subsubsection{Immutability}
% Functional programming encourages stateless functions which operate on immutable data structures and objects.
% An immutable object is an object that once initialized, cannot be modified.
% Immutability offers several benefits to software development in areas including concurrency and testing.
% Listing~\ref{lst:immutable} shows an example of manipulating an immutable list.
% \lstinputlisting[label={lst:immutable}, float, caption=Manipulating immutable list]{code/immutable.scala}
% Note that each \texttt{map} and \texttt{filter} operation creates a new copy of the list with the modified contents.
% The original list remains unchanged.
% Listing~\ref{lst:mutable} show the equivalent operation using a mutable list.
% \lstinputlisting[label={lst:mutable}, float, caption=Manipulating mutable list]{code/mutable.scala}
% Observe that listing~\ref{lst:immutable} is a single statement while listing~\ref{lst:mutable} contains multiple statements.

\subsubsection{SBT}\label{sec:sbt}
SBT\autocite{_sbt_????} is an interactive build tool used by many Scala projects.
SBT configuration files are written in \texttt{*.sbt} or \texttt{*.scala} files using Scala syntax and semantics.
Although SBT configuration files use plain Scala, they typically use coding patterns which are different from traditional Scala programs.
Listing~\ref{lst:sbt} is an example project definition in SBT.
\lstinputlisting[label={lst:sbt}, float, caption=SBT project definition]{code/sbt.scala}
Observe that the project is defined as a single statement and makes extensive use of symbolic infix operators.
Due to the nature of build configurations, argument lists can becomes unwieldy long and a single project statement can span over dozens or even hundreds of lines of code.


\subsection{scala.meta}\label{sec:scalameta}
Scala.meta\autocite{scala57:online} is a metaprogramming toolkit for Scala.
Before scala.meta, most metaprogramming facilities relied on Scala compiler internals.
This had several severe limitations such as too-eager desugaring resulting in loss of syntactic details from the original source code.
Scala.meta was designed to overcome these limitations and offer a more robust platform to develop metaprogramming tools for Scala.
Several key features of scala.meta have made it an invaluable companion in the development of scalafmt.
Most notably among these features are dialect agnostic syntax trees, syntax tree serialization, high-fidelity  parsing and algebraically typed tokens.

Scala.meta provides facilities to tokenize and parse a variety of different Scala dialects.
One such dialect is SBT configuration files, discussed in section~\ref{sec:sbt}.
SBT adds custom support for top-level statements in \texttt{*.sbt} files, a disallowed feature in regular Scala programs.
Then, to format SBT files requires either depending on the SBT parser or reimplementing its parsing logic.
To add insult to injury, top-level statements must be separated by a blank line if you use SBT version 0.13.6 or lower; a restriction that was lifted in SBT 0.13.7.
Listing~\ref{lst:dialect} shows how scala.meta dialects makes it trivial to accommodate this zoo of nuances.
\lstinputlisting[label={lst:dialect}, float, caption=Parsing different Scala dialects with scala.meta]{code/dialect.scala}
The result after parsing is a dialect agnostic scala.meta tree structure.

The structure of scala.meta trees can be serialized to a string to strip off all insignificant syntactic details.
Listing~\ref{lst:structure} shows how to serialize the tree structure of a simple hello world application.
\lstinputlisting[label={lst:structure}, float, caption=Serializing scala.meta trees]{code/structure.scala}
For example, observe that the comment has been stripped away.
As we discuss in section~\ref{sec:tooling}, this feature was instrumental in testing scalafmt.

Node types in scala.meta trees preserve absolute fidelity with the original source file.
This means we can obtain all syntactic details from a tree node such as whether a for comprehension uses parentheses or curly braces as delimiters, whitespace positions, comments and other syntactic trivia.
The Scala compiler is infamous for desugaring for-comprehensions into \texttt{map/withFilter/flatMap} applications during the parse phase.
This made it impossible to implement metaprogramming tasks such as code formatting.
High-fidelity parsing in scala.meta has been essential for scalafmt because we can't lose critical syntactic details such as whether for-comprehensions are used over \texttt{flatMap} method calls.

Tokens in scala.meta are strongly typed.
Traditional object-oriented libraries treat tokens as a single type with multiple methods such as \texttt{isComma/isFor} which returns true if a token instance is a comma or a \texttt{for} keyword.
However, scala.meta leverages algebraic data types in Scala to represent each different kind of token as a separate type.
This feature plays nicely with exhaustivity checking in the Scala pattern matcher and enabled design pattern for the \emph{Router} explained in section~\ref{sec:router}.


% \subsection{Metaprogramming with scala.meta}
% Scala.meta\autocite{_scala.meta_????} is a robust and portable metaprogramming toolkit for Scala.
% Scala 2.10 compiler introduced the first metaprogramming API for Scala via compile-time macros and runtime reflection.
% Library authors could take advantage of this new functionality to implement more advanced features in their libraries.
% However, the Scala compiler metaprogramming facilities exposed compiler internals to its users.

\subsection{Code formatting}
%
% In the literature, the term \emph{pretty-printer} is more commonly used than code formatter.
% This thesis uses code formatter the terms are used interchangeably and the author does not distinguish between the two.} 
% We emphasize formatters that support line length length limits.
% As future sections discuss, line length limits turn out to be a big challenge to implement.

Code formatting and pretty printing\footnote{
  This thesis uses the term \emph{code formatting} over \emph{pretty printing}.
  According to Hughes\autocite{hughes_design_1995}, pretty printing is a subset of code formatting where
  the former is only concerned with presenting data structures while the latter is
  concerned with the harder problem of formatting existing source code --- the main topic of this thesis.} has a long tradition.
In this chapter, we look at a variety of tools and algorithms that have been developed over the last 70 years.

\subsubsection{Natural language}
The science of displaying aesthetically pleasing text dates back as early as 1956\autocite{harris_keyboard_1956}.
The first efforts involved inserting carriage returns in natural language text.
Until that time, writers had been responsible for manually providing carriage returns in their documents before sending them off for printing.
The motivation behind automating this process was to ``save operating labor and reduce human error''.
Once type-setting became more commonplace, the methods for breaking lines of text got more sophisticated.

Knuth and Plass developed in 1981 a famous line breaking algorithm\autocite{knuth_breaking_1981} for \LaTeX{}, a popular typesetting program among scientific circles.
\LaTeX{} is the program that was used to generate this very document.
The line breaking problem was the same as in the 60s: how to optimally break a paragraph of text into lines so that the right margin is minimized.
The primitive approach is to greedily fit as many words on a line as possible.
However, such an approach can produce embarrassingly bad output in the worst case.
Knuth's algorithm uses dynamic programming to find an optimal layout with regards to a fit function that penalizes empty space on the right margin of the paragraph.
This algorithm remains a textbook example of an application of dynamic programming\autocite{dreyfus2002richard,kleinberg2006algorithm}.

\subsubsection{ALGOL 60}
Scowen\autocite{scowen_soapprogram_1971} developed SOAP in 1971, a code formatter for the programming language ALGOL 60.
The main motivation for SOAP was to make it ``easier for a programmer to examine and follow a program'' as well as to maintain a consistent coding style.
This motivation is still relevant in modern software development.
SOAP did provide a maximum line length limit.
However, SOAP would fail execution if the provided line length turned out to be too small.
With hardware from 1971, SOAP could format 600 lines of code per minute.

\subsubsection{LISP}\label{sec:lisp}
In 1973, Goldstein\autocite{goldstein_pretty-printing_1973} explored code formatting algorithms for LISP\autocite{mccarthy_recursive_1960} programs.
LISP is a family of programming languages and is famous for its parenthesized prefix notation.
Listing~\ref{lst:lisp} shows a program in LISP to calculate factorial numbers.
\lstinputlisting[label={lst:lisp}, float, caption=A LISP program]{code/lisp.lisp}
The simple syntax and extensive use of parentheses as delimiters makes LISP programs an excellent ground to study code formatters.

Goldstein presented a \emph{recursive re-predictor} algorithm in his paper.
The recursive re-predictor algorithm runs a top-down traversal on the abstract syntax tree of a LISP program.
While visiting each node, the algorithm tries to first obtain a \emph{linear-format},
i.e. fit remaining body of that node on a single line,
with a fallback to \emph{standard-format},
i.e. each child of that node is put on a separate line aligned by the first child.
Goldstein observes that this algorithm is practical despite the fact that its running time is exponential in the worst case.
Bill Gosper used the re-predictor algorithm to implement GRINDEF\autocite{_bill_????}, one of the first code formatters for LISP.

Goldstein's contributions extend beyond formatting algorithms.
Firstly, in his paper he studies how to format comments.
Secondly, he presents several different formatting layouts which can be configured by the users.
Both are relevant concerns for modern code formatters.

\subsubsection{Language agnostic}\label{sec:agnostic}
Derek C. Oppen pioneered the work on language agnostic code formatting in 1980\autocite{oppen_prettyprinting_1980}.
A language agnostic formatting algorithm can be used for a variety of programming languages instead of being
tied to a single language.
Users provide a preprocessor to integrate a particular programming language with the algorithm.
Oppen's algorithm runs in $O(n)$ time and uses $O(m)$ memory for an input program of length $n$ and maximum column width $m$.
Besides impressive performance results, Oppen claims that a key feature of his algorithm is its streaming nature; the algorithm prints formatted lines as soon as they are read from the input instead of waiting until the entire input stream has been read.
This feature is typically not a concern for modern code formatters.
Moreover, Oppen's algorithm shares a worrying limitation with SOAP: it cannot handle the case when the line length is insufficiently large.

Mark van der Brand presented a library in 1996 that generates a formatter given a context-free grammar\autocite{van_den_brand_generation_1996}.
% Brand correctly identifies that the development of code formatters requires a lot of effort.
Beyond the usual motivations for developing code formatters, Brand mentions that formatters ``relieve documentation writers from typesetting programs by hand''.
This focus on documentation is reflected by the fact that the generated formatter could produce both ASCII formatted code as well at \LaTeX{} markup.
Since comments are typically not included in a syntax tree, the presented algorithm has an elaborate scheme to infer the location of comments in the produced output.
Like Oppen's algorithm, this library requires the user to plug in a preprocessor to integrate a particular programming language into the Brand's library.
Unlike Oppen's algorithm, Brand does not consider line length limits in his algorithm.

John Hughes extended on Oppen's work on language agnostic formatting in terms of functional programming techniques\autocite{hughes_design_1995}.
Hughes presented a design of a \emph{pretty-printing} library that leverages combinators with algebraic properties to express formatting layouts.
Hughes claims that such a formal approach was invaluable when designing the pretty-printing library, which has seen wide use, including in the Glasgow Haskell compiler.
Wadler\autocite{wadler_prettier_2003} and Chitil\autocite{swierstra_linear_2009} extend on Hughes's and Oppen's work in term of performance and programming techniques.
However, this branch of work has been limited to printing data structures and not how to format existing source code.

\subsubsection{Go}
\gofmt\autocite{gofmt94:online} is a code formatter for the Go programming language, developed at Google.
\gofmt{} was released in the early days of Go in 2009 and is noteworthy for its heavy adoption by the Go programming community.
Official Go documentation\autocite{CodeR59:online} claims that almost all written Go code --- at Google and elsewhere --- is formatted with \gofmt{}.
Besides formatting, \gofmt{} is used to automatically migrate Go codebases from legacy versions to new source-incompatible releases.
However, \gofmt{} supports neither a maximum line length setting nor an opinionated setting.
Line breaks are preserved in the user's input.
For example, listing~\ref{lst:gofmt} shows a Go program that uses the same layout as the ``unformatted'' code from listing~\ref{lst:unformatted} in the introduction.
\lstinputlisting[label={lst:gofmt}, caption=Gofmt example input/output]{target/gofmt.go}
The output of running \gofmt{} through listing~\ref{lst:gofmt} is identical to the input.
This un-opinionated behavior may be considered desirable by many software developers.
However, in this thesis, we are concerned with maximum line lengths and opinionated settings.

\subsubsection{Scala}\label{sec:scalariform}
Scalariform\autocite{russell_scalariform_2010} was released in 2010 and is a widely used code formatter for Scala.
Like \gofmt{}, Scalariform does an excellent job of tidying common formatting errors.
Moreover, Scalariform supports a variety of configuration options.
Scalariform is also impressively fast, it can format large files with over 4.000 lines of code in under 250 milliseconds on a modern laptop.
However, Scalariform shares the same limitations with \gofmt{}: it lacks a line length and opinionated setting.

Firstly, the line length setting is necessary to implement many popular coding styles in the Scala community.
For example, the Spark\autocite{xin_spark_2015} and Scala.js\autocite{doeraene_scala.js_2015} coding styles have 100 character and 80 character column limits, respectively.
As we see in other code formatters, it is non-trivial to add a line length setting and doing so would require a significant redesign of Scalariform.

Secondly, the lack of an opinionated setting makes it impossible to enforce certain coding styles.
For example, the Scala.js coding style enforces \emph{bin-packing}, where arguments should be arranged compactly up to the column length limit.
Listings~\ref{lst:bin_pack} and~\ref{lst:non_bin_pack} shows an example of bin packing enabled and disabled, respectively.

\begin{minipage}{.45\textwidth}
  \lstinputlisting[label={lst:bin_pack}, caption=Bin-packing]{code/bin-pack.scala}
\end{minipage}
\hfil
\begin{minipage}{.45\textwidth}
  \lstinputlisting[label={lst:non_bin_pack}, caption=No bin-packing]{code/non-bin-pack.scala}
\end{minipage}

Since Scalariform preserves the line breaking decisions from the input,
Scalariform is unable to convert formatted code like in listing~\ref{lst:non_bin_pack} to the code in listing~\ref{lst:bin_pack}.

\subsubsection{C-family}
Daniel Jasper triggered a new trend in optimization based coded formatters with the release of \emph{ClangFormat}\autocite{jasper_clang-format_2014} in 2013.
ClangFormat is developed at Google and can format an impressive number of languages: C, C++, Java, JavaScript, Objective-C and Protobuf code.
Figure~\ref{fig:clang_format} shows the architecture of ClangFormat.
\begin{figure}[H]
  \centering
  \includegraphics[width=0.9\textwidth]{img/clang-format.png}
  \caption{ClangFormat architecture, source \autocite{jasper_clangformat_2013}}
  \label{fig:clang_format}
\end{figure}
The main components are the \emph{structural parser} and the \emph{layouter}.

ClangFormat employs a structural parser to split source code into a sequence of \emph{unwrapped-lines}.
An unwrapped line is a statement that should fit on a single line if given sufficient line length.
A key feature of unwrapped lines is that they should not influence other unwrapped lines.
The parser is lenient and parses even syntactically invalid code.
The parsed unwrapped lines are passed onto the layouter.

The ClangFormat layouter uses a novel approach to implement line wrapping.
Each line break is assigned a penalty according to several rules such as nesting and token type.
At each token, the layouter can choose to continue on the same line or break the line.
This forms an acyclic weighted directed graph with non-whitespace tokens representing vertices and splits (e.g., space, no space or line break) representing edges.
The first token of an unwrapped line is the root of the graph and the last token is the sink.
The layouter uses Dijkstra's\autocite{dijkstra_note_1959} shortest path algorithm to find the layout that has the lowest penalty.
To obtain good performance, the layouter uses several domain specific optimizations to minimize the search space.

Despite supporting several programming languages, ClangFormat does not leverage the language agnostic formatting techniques described section~\ref{sec:agnostic}.
Support for each language has been added as ad-hoc extensions to the ClangFormat parser and layouter.
ClangFormat supports a variety of configuration options, including 6 out-of-the-box styles based on coding styles from Google, LLVM and other well-known organizations.

Finally, ClangFormat is opinionated.
ClangFormat produces well-formatted output for even the most egregiously formatted input.
Listing~\ref{lst:clang_opinion} shows an offensively formatted C++ code snippet.
\lstinputlisting[label={lst:clang_opinion}, float, language=c, caption=Unformatted C++ code]{code/terrible.cpp}
Listing~\ref{lst:clang_opinion2} shows the same snippet after being formatted with ClangFormat.
\lstinputlisting[label={lst:clang_opinion2}, float, language=c, caption=ClangFormat formatted C++ code]{target/unterrible.cpp}
Observe that ClangFormat does not respect the (lack of) line breaking decisions in listing~\ref{lst:clang_opinion}.
This feature makes it possible to ensure that all code follows the same style guide, regardless of author.

\subsubsection{Dart}
Dartfmt\autocite{nystrom_dart_style_2014} was released in 2014 and follows the optimization based trend initiated by ClangFormat.
Dartfmt is a code formatter for the Dart programming language, developed at Google.
Like ClangFormat, dartfmt has a line length setting and is opinionated.
Bob Nystrom, the author of dartfmt, discusses the design of dartfmt in a blog post\autocite{nystrom_hardest_2015}.
In his post, Nystrom argues that the design of a code formatters is significantly complicated by a column limit setting.
The line wrapping algorithm in dartfmt employs a \emph{best-first search}\autocite{pearl_heuristics:_1984},
a minor variant of the shortest path search in ClangFormat.
As with ClangFormat, a range of domain-specific optimizations were required to make the search scale for real-world code.
Listing~\ref{lst:dead_end} shows an example of such an optimization: \emph{avoid dead ends}.
\lstinputlisting[label={lst:dead_end}, float, caption=Avoid dead ends]{code/dead-end.dart}
Line 4. in the snippet exceeds the 35 character column limit.
A plain best-first search would fruitlessly explore a lot of line breaking options inside the argument list of \texttt{firstCall}.
However, \texttt{firstCall} already fits on a line and there is no need to explore line breaks inside its argument list.
The dartfmt optimized search is able to eliminate such dead ends and quickly figure out to break before the \texttt{"long argument string"} literal.

\subsubsection{R}
The most recent addition to the optimization based formatting trend is \rfmt\autocite{yelland_new_2016},
a code formatter for the statistical programming environment \emph{R}.
The formatter was released in 2016 -- after the background work on this thesis started -- and like its forerunners is also developed at Google.
\rfmt{} makes an interesting contribution in that it combines the algebraic combinator approach from Hughes\autocite{hughes_design_1995} and the optimization based approach from \LaTeX{} and ClangFormat.

The algebraic combinator approach employed by rfmt makes it easy to express a variety of formatting layouts.
\rfmt{} uses 6 layout combinators or \emph{blocks} as they are called in the report.
% , to express all formatting rules for the R programming language.
The blocks are the following:
\begin{itemize}
  \item $TextBlock(txt)$: unbroken string literal.
  \item $LineBlock(b_1, b_2, \ldots, b_n)$: horizontal combination of blocks.
  \item $StackBlock(b_1, b_2, \ldots, b_n)$: vertical combination of blocks.
  \item $ChoiceBlock(b_1, b_2, \ldots, b_n)$: selection of a best block.
  \item $IndentBlock(n, l)$: indent block $b$ by $n$ spaces.
  \item $WrapBlock(b_1, b_2, \ldots, b_n)$:
    Fit as many blocks on each line as possible, break when the column limit is exceeded and align by the first character in $b_1$.
    % Similar to \emph{bin-packing} in other code formatters such as ClangFormat.
\end{itemize}
We'll use an example to demonstrate how these relatively few combinators allow an impressive amount of flexibility.
Listings~\ref{lst:rfmt1} and~\ref{lst:rfmt2} shows two different layouts to format an argument list.

\begin{minipage}{.45\textwidth}
  \lstinputlisting[label={lst:rfmt1}, caption=Line block]{code/line.scala}
\end{minipage}
\hfil
\begin{minipage}{.45\textwidth}
  \lstinputlisting[label={lst:rfmt2}, caption=Stack block]{code/stack.scala}
\end{minipage}

In this case, we prefer the line block from listing~\ref{lst:rfmt1} since it requires fewer lines.
However, our preference changes if the function name is longer as is shown in listings~\ref{lst:rfmt3} and~\ref{lst:rfmt4}.

\begin{minipage}{.45\textwidth}
  \lstinputlisting[label={lst:rfmt3}, caption=Line block]{code/line1.scala}
\end{minipage}
\hfil
\begin{minipage}{.45\textwidth}
  \lstinputlisting[label={lst:rfmt4}, caption=Stack block]{code/stack1.scala}
\end{minipage}

Here, we clearly prefer the stack block in listing~\ref{lst:rfmt4}.
Listing~\ref{lst:rfmt-example} shows how we use the 6 fundamental blocks in the \rfmt{} combinator algebra to express the choice between these two formatting layouts.
\begin{minipage}{.45\textwidth}
  \lstinputlisting[numbers=none, label={lst:rfmt-example},
                 caption=Formatting layout for argument lists]{code/rfmt-example.scala}
\end{minipage}
The variable $f$ denotes the function name and $a_1, ... a_m$ denotes the argument list.
Observe that listing~\ref{lst:rfmt-example} does not express how to find the optimal layout.

To find an optimal layout, \rfmt{} employs dynamic programming with a novel indexing scheme.
First, it is possible enumerate all layout combinations like the re-predictor algorithm does in section~\ref{sec:lisp}.
This leads to exponential growth which turns out to be a problem for some cases.
Dynamic programming alleviates exponential growth by allowing us to reuse partial solutions.
Instead of re-calculating the layout cost at each (starting column, block) pair, we store the result in an associative array keyed by the starting column.
However, it turns out that this can still be inefficient in terms of memory and speed\footnote{
  In fact, ClangFormat started with a similar approach, as explained in this\autocite{clang84:online} video recording, but then switched to Dijkstra's shortest path algorithm.
}.
To overcome this limitation, Yelland -- the \rfmt{} author -- presents an indexing scheme that makes it possible to extrapolate the layout cost even for missing keys.
We refer to the original paper\autocite{yelland_new_2016} for details.
This novel approach enables \rfmt{} to format even the most pathologically nested code in near instant time.

% TODO?
% \subsubsection{rustfmt} 


\section{Algorithms}\label{sec:algorithms}
This chapter describes how scalafmt formats Scala code.
We will see that scalafmt's design is inspired by ClangFormat and dartfmt.
However, we believe our design makes a valuable contribution in that it leverages functional programming principles to maximise code reuse and extensibility (TODO, really???).

\subsection{Design}
Figure~\ref{fig:architecture} shows a broad architectural overview of scalafmt.
\begin{figure}
  \centering
  \includegraphics[width=\textwidth]{img/architechture.pdf}
  \caption{Scalafmt architecture}
  \label{fig:architecture}
\end{figure}
First, scalafmt parses a source file using scala.meta.
Next, we feed a sequence of \emph{FormatToken} data types into a \emph{LineWrapper}.
The LineWrapper uses a \emph{Router} to construct a weighted directed graph and run a best-first search to find an optimal formatting layout for the whole file.
Finally, the LineWrapper feeds a sequence of \emph{Split} data types into the \emph{FormatWriter}, which constructs a new reformatted source file.
The following sections explain these data types and abstractions in detail.

\subsection{Data structures}
Scalafmt leverages a few carefully designed data structure to allow an implementation that emphasizes correctness and maintainability.

\subsubsection{FormatToken}
A \emph{FormatToken} is a pair of two non-whitespace tokens.
Listing~\ref{lst:format_token} shows the definition of the FormatToken data type.
\lstinputlisting[label={lst:format_token}, caption=FormatToken definition]{code/format_token.scala}
As shown in the architecture overview in figure~\ref{fig:architecture}, each token except the beginning and end of file tokens appear twice in the sequence of FormatTokens: once as the \texttt{left} member and once as the \texttt{right} member.
In a nutshell, the job of the LineWrapper is to convert each FormatToken into a \emph{Split}

\subsubsection{Decision}\label{sec:decision}
A Decision is a pair of a FormatToken and a sequence of Splits.
Listing~\ref{lst:decision} shows the definition of Decision.
\lstinputlisting[label={lst:decision}, caption=Decision definition]{code/decision.scala}
The \emph{splits} member represents the possible splits that the LineWrapper can choose for \emph{formatToken}.

\subsubsection{Policy}\label{sec:policy}
A \emph{Policy} is an enforced formatting layout over a region.
Listing~\ref{lst:policy} shows the definition of Policy.
\begin{minipage}{\linewidth}
  \lstinputlisting[label={lst:policy}, caption=Policy definition]{code/policy.scala}
\end{minipage}
A Policy is a partial function that should be applied to future Decisions up until the \emph{expire} token.
Policies easily compose using the Scala standard library \texttt{orElse} and \texttt{andThen} methods on PartialFunction\footnote{
  Fun fact. Careful eyes will observe that Policy is in fact a monoid with the empty partial function as identity and function composition as associative operator.}.
% The \texttt{isProhibitive} member annotates whether this policy can possibly eliminate Splits from future Decisions.
% This is necessary for an optimization explained in section~\ref{sec:dequeue}.
Policies enable a high-level way to express arbitrary formatting layouts over a region of code.

\subsubsection{Indent}\label{sec:indent}
An \emph{Indent} describes indentation over a region of code.
\lstinputlisting[label={lst:indent}, caption=Indent definition]{code/indent.scala}
Listing~\ref{lst:indent} shows the definition of Indent along with the algebraic data type \emph{Length}.
Length can either be \texttt{Num(n)} where $n$ represents an explicit number of spaces to indent by or \texttt{StateColumn} which is a placeholder for the number of spaces required to vertically align by the current column.
Indent is type parameterized by Length so that, at some point, we can replace StateColumn placeholders with Nums to obtain a concrete number.
For example, given a scala.meta tree \texttt{expr}, the definition \texttt{Indent(Num(2), expr.tokens.last, inclusive=true)}
increases the indentation level by 2 spaces up to and including the last token of \texttt{expr}.
The \texttt{inclusive} member is set to false when the indentation should expire before the expire token, for example in a block wrapped by curly braces, since the closing curly brace should not be indented by 2 spaces.
The StateColumn placeholder is required to allow memoization of Splits, which is critical for performance reason as explained in section~\ref{sec:router} on the \emph{Router}.

\subsubsection{Split}
A \emph{Split} represents a (possibly empty) whitespace character to be inserted between two non-whitespace tokens.
Listing~\ref{lst:split} shows the rather intricate definition of the Split data type\footnote{
  For clarity reasons, a few less important members have been removed from the actual Split definition.}.
\begin{minipage}{\linewidth}
  \lstinputlisting[label={lst:split}, caption=Split definition]{code/split.scala}
\end{minipage}
The Split data type went through several generations of design before reaching its current structure.
Each member serves an important role.
The most important member of the Split type is the \emph{modification}.
A modification must be one of \texttt{NoSplit}, \texttt{Space} and \texttt{Newline}.
The \emph{cost} member represents the penalty for choosing this split.
The \emph{optimalToken} member enables an optimization explained in section~\ref{sec:optimal}.
The \emph{line} member allows a powerful debugging technique explained in section~\ref{sec:router}.
The \emph{policy} and \emph{indents} members are explained in sections~\ref{sec:policy} and~\ref{sec:indent}, respectively.

\subsubsection{State}
A \emph{State} is a partial formatting solution during the best-first search.
Listing~\ref{lst:state} shows the definition of the State class and companion object.
\begin{minipage}{\linewidth}
  \lstinputlisting[label={lst:state}, caption=State definition]{code/state.scala}
\end{minipage}
Observe the similarity of State and Split.
A State contains various summaries calculated from the \texttt{splits} vector.
The summaries are necessary for performance reasons in the best-first search.
Observe that the indents are type parameterized by \texttt{Num}, meaning they only contain concrete indentations and no \texttt{StateColumn} indents.
The \texttt{indentation} member is the sum of all currently active indents and \texttt{column} represents how many characters have been consumed since the last newline.
The State class extends the \texttt{Ordered} trait to allow for efficient polling from a priority queue.
The \texttt{compare} method orders States firstly by their \texttt{totalCost} member, secondly by \texttt{splits.length} -- how many FormatTokens have been formatted -- and finally breaking ties by the \texttt{indentation}.
The method \texttt{nextState} calculates a penalty for characters that overflow the column limit and prepares
The method is implemented as efficiently as possible since the method is on a hot path in the best-first search.

\subsection{LineWrapper}
The LineWrapper is responsible for turning FormatTokens into Splits.
To accomplish this, the LineWrapper employs a \emph{Router} and a \emph{best-first search}.

\subsubsection{Router}\label{sec:router}
The Router's role is to produce a Decision given a FormatToken.
Figure~\ref{fig:router} shows all possible formatting layout for the small input \texttt{val x = y + z}.
\begin{figure}
  \centering
  \includegraphics[width=\textwidth]{target/router.pdf}
  \caption{Example graph produced by Router}
  \label{fig:router}
\end{figure}
In this figure, the Router is the planner that has chosen to open up multiple branches at \texttt{=} and \texttt{+} and only one branch for the remaining tokens.
This is no easy task since a FormatToken can be any pair of two tokens.
How do we go about implementing a Router?

The Router is implemented as one large pattern match on a FormatToken.
Listing~\ref{lst:match} shows how we can pattern match on a FormatToken and produce Splits.
\lstinputlisting[label={lst:match}, caption=Pattern matching on FormatToken]{code/match.scala}
The pattern \texttt{\_: `=`} matches a scala.meta token of type \texttt{`=`}.
The underscore \texttt{\_} ignores the underlying value.
Keyword is a super-class of all scala.meta keyword token types.
Now, a good observer will notice that this pattern match can quickly grow unwieldy long once you account for all of Scala's rich syntax.
How does this solution scale?
Also, once the match grows bigger how can we know from which case each Split origins?
It turns out that Scala's pattern matching and scala.meta's algebraically typed tokens are able to help us.

The Scala compiler can statically detect unreachable code.
If we add a case that is already covered higher up in the pattern match, the Scala compiler issues a warning.
For example, listing~\ref{lst:exhaustive} shows an example where the compiler issues a warning.
\begin{minipage}{\linewidth}
  \lstinputlisting[label={lst:exhaustive}, caption=Unreachable code]{code/exhaustive.scala}
\end{minipage}
Here, we accidentally match on a FormatToken with an \texttt{else} keyword on the right which will never match because we have a broader match on a Keyword higher up.
In this small example, the bug may seem obvious but once the Router grows bigger the compiler becomes unmissable.
However, this still leaves us with the second question of finding the origin of each Split.
Scala macros\autocite{burmako2013scala} and implicits\autocite{oliveira2010type} come to the rescue.

The source file line number of where a Split is instantiated is automatically attached on each Split.
Remember in listing~\ref{lst:split} that the Split case class had an implicit member of type \texttt{sourcecode.Line}.
Sourcecode\autocite{lihao91:online} is a tiny Scala library to extract source code metadata from your programs.
The library leverages Scala macros and implicits to unobtrusively surface useful information such as line number of call sites.
Listing~\ref{lst:sourcecode} shows how this works.
\lstinputlisting[label={lst:sourcecode}, caption=Extracting line number from call site]{code/sourcecode.scala}
When a \texttt{sourcecode.Line} not passed explicitly as an argument to Split's constructor, the Scala compiler will trigger its implicit search to fill the missing argument.
The \texttt{sourcecode.Line} companion contains an implicit macro that generates a Line instance from an extracted line number.
Take a moment to appreciate how these two advanced features of the Scala programming language enable a very powerful debugging technique.
The scalafmt router implementation contains 88 cases and spans over 1.000 lines of code.
The ability to trace the origin of each Split has been indispensable in the development of the Router.
% Once the Router can produce Decisions, we can run a best-first search to choose the optimal splits.

\subsubsection{Best-first search}
The Decisions from the Router produce a directed weighted graph, as demonstrated in figure~\ref{fig:router}.
To find the optimal formatting layout, our challenge is to find the cheapest path from the root node to a final node.
The best-first\autocite{pearl_heuristics:_1984} algorithm is an excellent fit for the task.

Best-first search is a graph search algorithm to efficiently traverse a directed weighted graph.
The objective is reach the final token and once we reach there, we terminate the search because we're guaranteed no other solution is better.
Algorithm~\ref{alg:bfsv1} shows a first attempt\footnote{
  We make heavy use of mutation since graph search algorithms typically don't lend themselves well to functional programming principles.
} to adapt a best-first search algorithm to the data structures and terminologies introduced so far.
\begin{algorithm}
  \caption{Scalafmt best-first search, first approach}\label{alg:bfsv1}
  \lstinputlisting[nolol]{code/bfsv1.scala}
\end{algorithm}
In the best case, the search always chooses the cheapest splits and the algorithm runs in linear time.
Observe that the router is responsible for providing well-behaved splits so that we never hit on the error condition after the while loop.
Excellent, does that mean the search is complete?
Absolutely not, this implementation contains several serious performance issues.

Algorithm~\ref{alg:bfsv1} is exponential in the worst case.
For example, listing~\ref{lst:exponential} shows a tiny input that triggers the search to explore over 8 million states.
\lstinputlisting[label={lst:exponential}, caption=Exponential running time]{code/exponential.scala}
Even if we could visit 1 state per microsecond\footnote{
  Benchmarks reveal the best-first search visits on average one state per 10 microseconds
} the search will take almost 1 second to complete.
This is unacceptable performance to format only 2 lines of code.
Of course, we could special-case long comments, but that would only provide us a temporary solution.
Instead, like with ClangFormat and dartfmt, we apply several domain specific optimizations.
In the following section, we discuss the optimizations that have shown to work well for scalafmt.

\subsection{Optimizations}\label{sec:optimizations}
This section explains the most important domain-specific optimizations that were required to get acceptable performance for scalafmt.
We will see that some optimizations are quite ad-hoc and require some creative workarounds.

\subsubsection{dequeueOnNewStatements}\label{sec:dequeue}
Once the search reaches the beginning of a new statement, empty the priority queue.
Observe that the formatting layout for each statement is independent from the formatting layout of the previous statement.
Consider listing~\ref{lst:dequeue}.
\lstinputlisting[label={lst:dequeue}, caption=Two independent statements]{code/dequeue.scala}
Both statements exceed the column limit, which means that the search must back-track to some extent.
However, once the search reaches \texttt{statement2} we have already found an optimal formatting layout for \texttt{statement1}.
When we start backtracking in \texttt{statement2}, there is no need to explore alternative formatting layouts for \texttt{statement1}.
Instead, we can safely empty the search queue once we reach the \texttt{statement2} token.

The \texttt{dequeueOnNewStatements} optimization is implemented by extending algorithm~\ref{alg:bfsv1} with an if statement.
Algorithm~\ref{alg:dequeue} shows a rough sketch of how this is done.
\begin{algorithm}
\caption{dequeueOnNewStatements optimization}\label{alg:dequeue}
  \lstinputlisting[nolol]{code/dequeue-alg.scala}
\end{algorithm}
With an empty queue, we ensure the search backtracks only as far back as is needed.
The \texttt{statementStarts} variable contains all tokens that begin a new statement.
To collect those tokens, we traverse the syntax tree of the input source file and select the first tokens of each statement of a block, each case in a partial function, enumerator in a for comprehension and so forth.
The actual implementation is quite elaborate and is left out of this thesis for clarity reasons.
Unfortunately, our optimization has one small problem.

Algorithm~\ref{alg:dequeue} may dequeue too eagerly inside nested scopes, leading the search to hit the error condition.
Listing~\ref{lst:noopt} shows an example where this happens.
\lstinputlisting[label={lst:noopt}, caption=Overeager dequeueOnNewStatements]{code/noopt.scala}
Remember that each case of a partial function starts a new statement.
The \texttt{dequeueOnNewStatements} optimization will dequeue the queue on the first state that reaches the \texttt{case} token.
In this example, the first state to reach the \texttt{case} token will have a strict Policy that disallows newlines up until the closing parenthesis.
However, we must insert a newline after the comment.
This causes the search to terminate too early and reach the error condition.
By inspecting where this problem occurred, we came up with a simple rule to identify regions where the \texttt{dequeueOnNewStatements} optimization should be disabled.
The simple rule is to never run \texttt{dequeueOnNewStatements} inside a pair of parentheses.
In section~\ref{sec:tooling}, we discuss techniques we used to be confident that this rule indeed works as intended.
In the following section (\ref{sec:recurseOnBlocks}) we explain the \texttt{recurseOnBlocks} optimization, which allows us to reenable \texttt{dequeueOnNewStatements} for selected regions inside parentheses.

\subsubsection{recurseOnBlocks}~\label{sec:recurseOnBlocks}
If the \texttt{dequeueOnNewStatements} optimization is disabled and we start a new block delimited by curly braces, recursively run the best-first search inside the block.
The intuition here is that by recursively running the best-first search, we keep the priority queue small at each layer of recursion.
This allows us to run aggressive optimizations such as \texttt{dequeueOnNewStatements}.

The \texttt{recurseOnBlocks} optimization enables scalafmt to handle idiomatic Scala code where large bodies of higher order functions and blocks are passed around as arguments.
Remember from section~\ref{sec:scala} that  Scala makes it syntactically convenient to in higher order functions as arguments to other functions.
Listing~\ref{lst:recurse} shows an example where this happens and we trigger the \texttt{recurseOnBlocks} optimization.
\lstinputlisting[label={lst:recurse}, caption=recurseOnBlocks example]{code/recurse.scala}
The \texttt{dequeueOnNewStatements} optimization is disabled inside argument list.
The priority queue grows out bounds because the higher order function can have an arbitrary number of statements.

To implement the \texttt{recurseOnBlocks} optimization, we add an extension to algorithm~\ref{alg:bfsv1}.
Algorithm~\ref{alg:recurse} shows a rough sketch of how \texttt{recurseOnBlocks} is implemented.
\begin{algorithm}
\caption{recurseOnBlocks optimization}\label{alg:recurse}
  \lstinputlisting[nolol]{code/recurse-alg.scala}
\end{algorithm}
We change the signature to accept a starting State and token where we stop the search.
Observe that we guard against infinite recursion by not making a recursive call on \texttt{start.formatToken}.
With \texttt{recurseOnBlocks} and \texttt{dequeueOnNewStatements}, we have solved most problems caused by independent statements affecting the formatting layouts of each other.
Next, we leverage recursion again to help the search queue stay small.

\subsubsection{OptimalToken}\label{sec:optimal}
An OptimalToken is a hint from a Split to the best-first search that enables the search to early eliminate competing Splits.
Recall from listing~\ref{lst:split} that a Split has an \texttt{optimalToken} member.
Listing~\ref{lst:optimalToken} shows the definition of OptimalToken.
\lstinputlisting[label={lst:optimalToken}, float, caption=OptimalToken definition]{code/optimalToken.scala}
When the best-first search encounters a Split with a defined OptimalToken,
the best-first search makes an attempt to reach that token with a budget of 0 cost.
If successful, the search can eliminate the competing Splits.
If unsuccessful and the \texttt{killOnFail} member is true, the best-first search eliminates the Split.
Otherwise, the best-first search continues as usual.

By eliminating competing branches, we drastically minimize the search space.
Listing~\ref{lst:optimal} shows an example where the OptimalToken can be applied.
\lstinputlisting[label={lst:optimal}, float, caption=OptimalToken example]{code/optimal.scala}
Scalafmt supports 4 different ways to format call-site function applications.
This means that there will be $4^N$ number of open branches when the search reaches \texttt{UserObject} $N$.
To overcome this issue, we define an OptimalToken at the closing parenthesis.
The best-first search successfully fits the argument list of each \texttt{UserObject} on a single line, and eliminates the 3 other competing branches.
This makes the search run in linear time as opposed to exponential.

To implement the \texttt{OptimalToken} optimization, we add an extension to algorithm~\ref{alg:recurse}.
Algorithm~\ref{alg:optimal} sketches how the extension works.
\begin{algorithm}
  \caption{OptimalToken optimization}\label{alg:optimal}
  \lstinputlisting[nolol]{code/optimalToken-alg.scala}
\end{algorithm}
The \texttt{bestFirstSearch} method has a new \texttt{maxCost} parameter, which is the highest cost that a new splits can have.
Next, if a Split has defined an \texttt{OptimalToken} we make an attempt to format up to that token.
If successful, we update \texttt{optimalFound} variable to eliminate other Splits from being added to the queue.
If unsuccessful and \texttt{killOnFail} is true, we eliminate the Split that defined the OptimalToken.
A straightforward extension to this algorithm would be to add a \texttt{maxCost} member to the \texttt{OptimalToken} definition from listing~\ref{lst:optimalToken}.
However, this has not been necessary for scalafmt.

\subsubsection{pruneSlowStates}
The pruneSlowStates is a optimization that eliminates states that progress slowly.
A state progresses slowly if it visits a token after other equally or less expensive states.
The insight is that if two equally expensive states visit the same token, the first state to visits that token typically produces a better formatting layout.

By eliminating slow states, we obtain a better formatting output in addition to minimizing the search space.
Listing~\ref{lst:slow} shows two formatting solutions that the Router has labelled as equally expensive.
However, the fast solution is explored first by the best-first search and, hence, we call it \emph{faster}.
\lstinputlisting[label={lst:slow}, caption=Slow states]{code/slow.scala}
The \texttt{pruneSlowStates} ensures that fast solutions are prioritized over slow solutions.
Of course, the Router could have assigned different costs to the line break after \texttt{g +} and \texttt{c +}.
However, our experience was that such as solution would introduce unnecessary complexity in the design of the Router.
Instead, the \texttt{pruneSlowStates} is able to take care of eliminating the slow state.

The \texttt{pruneSlowStates} is implemented as a extension to algorithm~\ref{alg:optimal}.
Algorithm~\ref{alg:slow} shows a rough sketch of how the extension works.
\begin{algorithm}[H]
  \caption{pruneSlowStates optimization}\label{alg:slow}
  \lstinputlisting[nolol]{code/slow-alg.scala}
\end{algorithm}
Observe that this algorithm is transparent to the Router.
No special annotations are required from Splits.

\subsubsection{escapeInPathologicalCases}
Alas, despite our best efforts to keep the search space small, some inputs can still trigger exponential running times.
The \texttt{escapeInPathologicalCases} optimization is our last resort to handle such pathological inputs.
How do we detect that the search has encountered such a challenging input?

We detect the search space is growing out of bounds by tallying the number of visits per token.
If we visit the same token $N$ times, we can estimate the current branching factor to be around $log_2(N)$.
In scalafmt, we tune $N$ to be 256 so that the best-first search can split into two or more paths for up to 8 tokens.
When a token has been visited more than 256 times, we trigger the \texttt{escapeInPathologicalCases} optimization.
In the following paragraphs section, we present two alternative fallback strategies: \emph{leave unformatted} and \emph{best-effort}.

The simplest and most obvious fallback strategy is to leave the pathologically nested code unformatted.
This can be implemented by backtracing to the first token of the current statement and then reproduce the formatting input up to the last token of said statement.
This method is guaranteed to run linearly to the size of the input.
The responsibility is left to the software developer to a manually format her code, removing all the benefits of code formatting.
However, in some cases the software developer may prefer to get a decent yet suboptimal formatting output.

The best-effort fallback strategy applies heuristics to give a decent but suboptimal formatting output.
When a token is visited for the 256th time, we select two candidate states from the search queue and eliminate all other states.
The first state is the fastest state --- the state that has reached furthest into the token stream --- that is not bound a prohibitive single line policy.
Single line policies are policies that eliminate newline Splits.
The second state is the current state --- the slow state that visited the token for the 256th time.
The intuition is that the fast state has good formatting output so far but for is stuck on a challenging token for some reason.
The slow may paid a hefty penalty early causing it to move slowly but maybe the early penalty will yield a better output in the end.
Algorithm~\ref{alg:best-effort} shows an example of the best-effort strategy can be implemented as an extension to algorithm~\ref{alg:bfsv1}.
\begin{algorithm}
  \caption{best-effort fallback strategy}\label{alg:best-effort}
  \lstinputlisting[nolol]{code/best-effort.scala}
\end{algorithm}
The \texttt{isSafe} method on State returns true if the state contains prohibitive policies, derived from annotated metadata in Splits from the Router.
Observe that this algorithm will reapply the best-effort fallback until the search reaches the final token.
In scalafmt, we bound the number of this can happen with a final fallback to the unformatted strategy.

The unformatted and best-effort fallback strategies offer different trade-offs.
The unformatted strategy works well in a scenario where a software developer is available to manually fix formatting errors.
The best-effort strategy works well on computer generated code where just a modicum of formatting still greatly aid the legibility of the code.
Unfortunately, as we discuss in section~\ref{sec:testing}, we struggled to guarantee idempotency using the best-effort strategy.
This limitation renders the best-effort strategy useless in environments where code formatters are used to enforce a consistent coding style across a codebase.


\subsection{FormatWriter}
Recall from figure~\ref{fig:architecture}, the FormatWriter receives splits from the best-first search and produces the final output presented to the user.
In addition to reifying Splits, the FormatWriter runs three post-processing steps: \emph{docstring formatting}, \emph{vertical alignment} and \emph{stripMargin alignment}.

\subsubsection{Docstring formatting}\label{sec:docstring}
Docstrings are used by software developers to document a specific part of code.
Like in Java, docstrings in Scala start with the \texttt{/**} pragma and end with \texttt{*/}.
However, unlike in Java, the Scala community is split on whether new lines inside docstrings should align by the first or the second asterisk.
The official Scala Style Guide\autocite{Scala80:online} dictates that new lines should align by the second asterisk while the Java tradition is to align by the first asterisk.
The Scala.js\autocite{doeraene_scala.js_2015} and Spark\autocite{xin_spark_2015} style guides follow the Java convention.
To accommodate all needs, scalafmt allows the user to choose either style.
To enforce that the asterisks are aligned according to the user's preferences,
the FormatWriter rewrites docstring tokens.
This is implemented with simple regular expressions.

\subsubsection{stripMargin alignment}
The Scala standard library adds a \texttt{stripMargin} extension method on strings.
The method helps Scala developers write multiline interpolated and regular string literals.
Listing~\ref{lst:stripMargin} shows an example usage of the \texttt{stripMargin} method.
\lstinputlisting[label={lst:stripMargin}, caption=stripMargin example]{code/stripMargin.scala}
After calling the method, the indentation and \texttt{|} character on line 3 are conveniently removed.
However, the hard-fought indentation on the pipe can easily be lost when the string is moved up or down a scope during refactoring.
Scalafmt can automatically fix this issue.
In the FormatWriter, scalafmt rewrites string literals to automatically align the \texttt{|} characters with the opening triple quotes \texttt{"""}.
This setting is disabled by default since scalafmt requires semantic information to confidently determine if the \texttt{stripMargin} invocation calls the standard library method or a user-defined method.

\subsubsection{Vertical alignment}
It turns out that vertical alignmyent is incredibly popular in the Scala programming community.
Vertical alignment is a formatting convention where similar tokens from different lines align together on the same column.
Listing~\ref{lst:align} shows an example of vertical alignment.
\lstinputlisting[label={lst:align}, caption=Vertical alignment example]{code/align.scala}
The lack of vertical alignment in the initial release of scalafmt was a hindrance for user-adoption.
Configurable vertical alignment was added to the 0.2 release.

Vertical alignment is implemented in the FormatWriter as an extension to the reification of Splits.
Instead of reifying a Space into a single space literal,
the FormatWriter builds a \texttt{tokenAligns: Map[Split, Int]} which specifies the number of additional spaces to add for each Space.
Algorithm~\ref{alg:align} shows a simplification of how the actual algorithm\footnote{
  The actual implementation is 130 lines of code, including helper methods. See
\href{https://git.io/voIG4}{https://git.io/voIG4}}
constructs the \texttt{tokenAligns} map.
\begin{algorithm}
  \caption{Vertical alignment, simplified algorithm}\label{alg:align}
  \lstinputlisting[nolol]{code/align-alg.scala}
\end{algorithm}
The running time of this algorithm is linear to the number of tokens.
In a nutshell, the algorithm builds blocks of lines that can be vertically aligned.
Blocks are separated by blank lines or mismatching candidates in subsequent lines.
As demonstrated in listing~\ref{lst:align}, key features of this algorithm include that users can configure arbitrary symbols to align by and lines can contain multiple columns of vertically aligned tokens.
This algorithm has two main limitations.
Firstly, all lines in a block must have an equal number of matching columns.
Secondly, the algorithm does not infer vertical alignment forcing the user to explicitly configure which tokens should align.

\subsubsection{Conclusion}
This section introduced the data structures and algorithms that scalafmt uses to implement line wrapping and vertical alignment.
Similar to dartfmt and ClangFormat, scalafmt models line wrapping like a graph search problem.
Non-whitespace tokens represent nodes and each potential split forms a weighted edge to the next token, where the associated cost represents the appeal of the formatting layout yielded by the split.
However, scalafmts implementation deviates quickly from there by introducing the concepts of Splits, Policies and Router.
We believe scalafmts use of partial functions in the Policy data type to enforce formatting layouts follows idiomatic Scala conventions, making it approachable for Scala developers to maintain and extend.
Similarly, we believe that translating dartfmts concepts of Rules, Spans and Chunks or ClangFormats concept of unwrapped lines to Scala would come at the price of less idiomatic Scala code.
Nevertheless, these concepts are different means to the same end.
We leave it to the judgement of the reader to assess which concepts are more powerful or intuitive to understand.

\section{Tooling}\label{sec:tooling}
This chapter describes the tools that we developed while designing an implementing algorithms for scalafmt.
These tools were indispensable in giving us confidence that our algorithms worked as intended.

\subsection{Heatmaps}
Section~\ref{sec:optimizations} introduces several extensions to algorithm~\ref{alg:bfsv1} that were required go get acceptable performance.
In general, the extensions involved eliminating search states.
To identify code patterns that triggered excessive search growth, we developed heatmaps.

Heatmaps are a visualization that displays which code regions are most frequently visited in the best-first search.
Figure~\ref{fig:heatmap} shows an example heatmap.
\begin{figure}
  \centering
  \includegraphics[width=\textwidth]{img/heatmap.png}
  \caption{Example heatmap with 5.121 visisted states}
  \label{fig:heatmap}
\end{figure}
The intensity of the red color indicates how often a particular token was visited.
A token highlighted by the lightest shade of red was visited twice while a token highlighted by the darkest shade of red was visited over 256 times.
This figure demonstrates several of the optimizations discussed in section~\ref{sec:optimizations}.
Firstly, thanks to the \texttt{dequeueOnNewStatements} optimization, the background is plain white up to the \texttt{Seq}.
The \texttt{Seq} gets visited twice, once when there's a space after the \texttt{=} and once when there's a newline.
Secondly, due to the OptimalToken optimization, when the search gets into trouble it backtracks to the tuple \texttt{(0, 0)} instead of the \texttt{Seq[((Int, Int), Matri)]} type signature.
Finally, because of the strategically placed comment at the end that exceeds the column limit, the search space grows out of bounds on the fourth argument triggering the \texttt{escapeInPathologicalCases} best-effort fallback.
Without heatmaps, it would be a much greater challenge to get these insights.
However, these heatmaps gave us limited insights in how our changes affected the search space in the best-first search.

We developed an extension to heatmaps that allows us to visually compare the difference in search space between two versions of scalafmt.
Figure~\ref{fig:heatmap2} shows an example of such a diff report.
\begin{figure}
  \centering
  \includegraphics[width=\textwidth]{img/heatmap2.png}
  \caption{Example diff heatmap}
  \label{fig:heatmap2}
\end{figure}
The green background indicates that the new version of scalafmt makes fewer visits to those regions.
Observe that the \texttt{>} operator has a background with a light shade of red.
This means that the operator was visisted more often in the new scalafmt version.
A price well worth paying considering the overall shrink in search space.
To produce diff heatmaps, we first persist to a database the statistics report needed to generate a single heatmap after each test run.
Then, we generate the diff heatmap by fetching two reports and calculating the difference in visits per token.
If the difference is negative for a particular token --- meaning we visited said token fewer times --- the background is highlighted green, otherwise red.






\subsection{Traceability}\label{sec:line}
\subsection{Configuration}
\subsubsection{maxColumn}
\subsubsection{binPacking}
\subsubsection{vertical alignment}
\subsection{Testing}~\label{sec:testing}
\subsection{Unit tests}
\subsection{Property based tests}
\subsubsection{AST Integrity}
\subsubsection{Idempotency}
\subsection{Regressions tests}

\section{Evaluation}
Code formatting is inherently a subjective topic.
This introduces a challenge when evaluating a code formatter.
In this chapter, we will present measurements that we believe show the success of scalafmt.
We do not measure how well software developers perceive scalafmt formatted code.
Instead, we will focus on \emph{performance benchmarks} and \emph{user adoption}.

\subsection{Performance benchmarks}
In this chapter, we describe our test methodology, key metrics and the performance results.

\subsubsection{Test methodology}
We run micro benchmarks to get insight into how scalafmt performs in an interactive software developer workflow where scalafmt runs on file save.
We select four files of varying sizes: small ($\sim$ 50 LOC), medium ($\sim$300 LOC), large ($\sim$1.000 LOC) and extra large ($\sim$4.500 LOC).
We use the OpenJDK Java Microbenchmark Harness (JMH)\autocite{OpenJ38:online} to measure the time it takes format the files.
JMH can be used to build, run and analyze benchmarks written in languages targeting the JVM, including Scala.
The sbt-jmh\autocite{ktoso84:online} plugin makes it easy to integrate JMH with a Scala project.
For accurate measurements, we run 10 warmup iterations followed by 10 measured iterations.
We compare the running time with Scalariform, an alternative code formatter for Scala discussed in section~\ref{sec:scalariform}.
The micro benchmark is single threaded.

We run macro benchmarks to get insights into how scalafmt performs in a continuous integration setup.
We format a sample 9.423 source files from seven open source libraries: intellij-scala, akka, spark, scalding, fastparse, scala-js and goose.
The whole sample contains a total of 1.219.235 lines of code.
Table~\ref{tab:macro_sample} shows the distribution of file sizes.
\begin{table}
  \centering
  \csvautobooklongtable{data/file_sizes.csv}
  \caption{Lines of code per file in macro benchmark}~\label{tab:macro_sample}
\end{table}
Observe that over 90 percent of all files contain under 300 lines of code.
Less than one percent of all files contain over 1.000 lines of code.
For accurate measurements, we run three iterations of the macro benchmark.
As with the micro benchmark, we compare the running time with Scalariform.
The macro benchmark is multi-threaded and takes advantage of all cores on the underlying hardware.

Our hardware is a Macbook Pro (Retina, 15-inch, Mid 2014) laptop with a quad-core 2.5 GHz Intel Core i7 processor, 256 KB L2 cache per core and 6 MB shared L3 cache.
The laptop has 16 GB 1600 MHz DDR3 memory.
The operating system is OS X El Capitan 10.11.5.
We run the benchmarks from the scalafmt commit id \href{https://github.com/olafurpg/scalafmt/tree/aff5e794dae4787b08243f8abb87a3ca4d907e40}{aff5e794} compiled against Scala 2.11.7, running on JVM version 8, update 91.

\subsubsection{Results}


\subsection{Adoption}\label{sec:adoption}
TODO




\section{Future work}
Scalafmt is not free from issues.
The main areas for future work regard how to more soundly produce an optimal formatting layout and how to obtain better performance in an interactive developer workflow.

We experienced a long tail of problems while getting the best-first search to reach the last token for a lot of Scala code that we found in the wild.
If the search cannot reach the last token, scalafmt cannot format the file.
Although the concept of policies does give a lot of flexibility to concisely express different formatting layouts, our experience is that it can be easy to create overly strict policies that eliminate all active search states.
It is worth to explore more principal approaches on how to define formatting layouts so that we can better guarantee that the search completes successfully every time.
We believe \rfmt{}s approach of combining the convenience of combinators with algebraic properties and optimized dynamic programming for excellent performance opens an interesting venue to solve this problem.

Incremental formatting opens an opportunity to get ludicrous formatting speed in an interactive developer workflow.
Incremental formatting works like incremental compilation in the sense that instead of formatting a whole file on every invocation, the formatter reuses output from previous invocations and only reformat the lines that have changed.
We believe that incremental formatting could cut down the formatting time for a large source file by several orders of magnitude, for example from 2s to 20ms.

\section{Conclusion}
We set out to implement a code formatter for the Scala programming language that supports several important features: an opinionated setting, a maximum line length setting, vertical alignment and fast performance.
In this thesis, we have presented data structures and algorithms that we used to develop scalafmt, a Scala code formatter that supports the first three required features and goes far towards achieving good performance.
Benchmarks reveal that scalafmt is 6x slower than Scalariform, an alternative Scala code formatter.
However, the quick user adoption of scalafmt indicates that scalafmt's current performance is already acceptable for many software developers.
We have reason to believe there is still plenty of room to improve on scalafmt's performance in order to meet the needs of the most performance demanding users.

% \section{Related work}
% \subsection{Combinator based}
% \begin{enumerate}
%   \item Houghes 1995
%   \item Wadler 1999
% \end{enumerate}
% \subsection{Optimization-oriented}
% \begin{enumerate}
%   \item clang-format Dijkstra's 2010
%   \item dartfmt Best-first search 2014
%   \item rfmt 2015
% \end{enumerate}
% \begin{enumerate}
%   \item Optimal line breaking
%   \item Oppen
% \end{enumerate}

\pagebreak
\printbibliography{}

\end{document}


